\documentclass[11pt]{article}

% \providecommand{\answerDirectory}{solutions}

../../macros/macros.tex

\inputAnswer{problem_0_answer}

\doHwHeader{Homework \BLUE{5}}

%%%%% EDIT THE FILES NAMED problem_X_answer.tex.  DON'T EDIT THIS FILE

%%%%%%%%%%%%%%%%%%%%%%%%%%%%%%%%%%%%%%%%%%%%%%%%%%%%%%%%%%%% 
%%%%%%%%%%%%%%%%%%%%%%%%%%%%%%%%%%%%%%%%%%%%%%%%%%%%%%%%%%%% 
%%%%%%%%%%%%%%%%%%%%%%%%%%%%%%%%%%%%%%%%%%%%%%%%%%%%%%%%%%%% 
%%%%%%%%%%%%%%%%%%%%%%%%%%%%%%%%%%%%%%%%%%%%%%%%%%%%%%%%%%%% 

\begin{document}


%%%%% EDIT THE FILES NAMED problem_X_answer.tex.  DON'T EDIT THIS FILE 

\paragraph{General instructions.}
% First read \BLUE{Lecture Note 5}, then
Answer the problems given in the rest of this document.
Edit this LaTeX template to add your answers, as described in the file named \texttt{00\_README.txt},
then compile it to produce \BLUE{\texttt{homework\_5.pdf}}\ifGradescope{} for submission to gradescope\fi.

\ifGradescope
For this homework, when you submit to gradescope you don't need to identify the pages
for each problem.  We're using gradescope's ``fixed-length'' grading mode.
To make it work, your answer for each part should fit in the allotted space,
so that each subsequent answer is on the expected page in the expected location.
Your PDF should have
\BLUE{9}
pages.  
\fi

\paragraph{Don't know?  Say so.}
If you don't know the answer to a problem or any part of the problem,
we encourage you to simply write \emph{``I don't know''} instead of giving an answer.
(If you write ``I don't know'' you can also add any questions you might have.)

\paragraph{Gentle reminder: homework is for learning.}
As noted in Lecture Note 1, engaging fully with the homeworks is the best way to learn the material.
\emph{The primary role of homeworks is to help you learn the material,
  and the primary role of feedback on homeworks is to help you improve your ideas.}

Please try to explain your ideas as clearly as you can.
The clearer your explanations are, the more useful the feedback on them can be.

\paragraph{Collaboration and citation policy.}
If you are doing this homework as part of a course, follow the course's collaboration and citation policy.
In any case, at the end of each answer, list the people you collaborated with,
and cite every source that your answer uses ideas from
(other than the lecture notes and the sources listed in the lecture notes).
If there are none, write ``none''.

%%%%% EDIT THE FILES NAMED problem_X_answer.tex.  DON'T EDIT THIS FILE 

\vfill
\noindent{\footnotesize\copyrightNotice}

%%% Local Variables:
%%% mode: latex
%%% TeX-master: "homework_2"
%%% End:


%%%%%%%%%%%%%%%%%%%%%%%%%%%%%%%%%%%%%%%%%%%%%%%%%%%%%%%%%%%% 
%%%%%%%%%%%%%%%%%%%%%%%%%%%%%%%%%%%%%%%%%%%%%%%%%%%%%%%%%%%% 
%%%%%%%%%%%%%%%%%%%%%%%%%%%%%%%%%%%%%%%%%%%%%%%%%%%%%%%%%%%% 
%%%%%%%%%%%%%%%%%%%%%%%%%%%%%%%%%%%%%%%%%%%%%%%%%%%%%%%%%%%% 

%%%%% EDIT THE FILES NAMED problem_X_answer.tex.  DON'T EDIT THIS FILE 

%%%%%%%%%%%%%%%% PROBLEM 1  %%%%%%%%%%%%%%%%%

% \setcounter{problem}{1}

\newpage
\problem \emph{(Greedy algorithm for \prob{Puncturing Intervals}; see Exercise 5.14)}

  Recall the \prob{Puncturing Intervals} problem from Homework 4: 

  \begin{mylist}
  \item[\bf input:] Set $I=\{J_1, J_2, \ldots, J_n\}$ of intervals, where $J_i = [s_i, e_i]$.
    We assume $e_1 \le e_2 \le \cdots \le e_n$.
  \item[\bf output:] A minimum-size finite set $P=\{p_1, p_2, \ldots, p_k\}$
    of points such that $J_i \cap P \ne \emptyset$ for each $J_i\in I$.
  \end{mylist}

  Consider the following rule for picking a single point $p_1$ that is guaranteed to be in some
  correct solution:

  \smallskip 

  \emph{1. Let $p_1$ be the end-time $e_1$ of the earliest-ending interval $J_1$.}

  \smallskip

  Consider the following two lemmas:

  \begin{lemma}\label{lemma: pi greedy choice}
    Let $I$ be any non-empty instance of \prob{Puncturing Intervals}.
    Let $e_1$ be the end-time of the earliest-ending interval $J_1$.
    Then $I$ has a correct solution that contains $e_1$.
  \end{lemma}

  \begin{lemma}\label{lemma: pi recurse}
    Let $I$ be any non-empty instance of \prob{Puncturing Intervals}.
    Let $e_1$ be the end-time of the earliest-ending interval $J_1$.
    Let $D$ contain the intervals in $I$ that don't contain $e_1$.
    Let $R$ be any optimal solution for $D$ (as a \prob{Puncturing Intervals} instance).
    Then $\{e_1\}\cup R$ must be a correct solution for $I$.
  \end{lemma}

  \begin{enumerate}[label=(\alph*)]
  \item Give a long-form proof of Lemma~\ref{lemma: pi greedy choice}.
    As much as possible, your proof should mimic the structure of the proof of Lemma 5.4 (for \prob{Activity Selection})
    in Lecture Note 5.

  \item Give a long-form proof of Lemma~\ref{lemma: pi recurse}.
    As much as possible, your proof should mimic the structure of the proof of Lemma 5.5 (for \prob{Activity Selection})
    in Lecture Note 5.

  \item Following the lemmas, 
    define a correct recursive algorithm \textsf{rpuncture} for \prob{Puncturing Intervals}.
    As much as possible, your algorithm should mimic the structure
    of the \textsf{rselect} algorithm for \prob{Activity Selection}
    in Section 5.3.
    Define the algorithm precisely in words, then give pseudo-code for it.

  \item Prove Theorem~\ref{thm: pi correctness}, below, for your algorithm.
    As much as possible, your proof should mimic the structure
    of the proof of Theorem 5.2 (for \textsf{rselect}) in Lecture Note 5.
    But you must give your proof in long form, following the template given here.
\end{enumerate}

\begin{theorem}\label{thm: pi correctness}
  For any instance $I=\{J_1, J_2, \ldots, J_n\}$ of \prob{Puncturing Intervals},
  \textsf{rpuncture}${(I)}$ returns a correct solution for $I$.
\end{theorem}

\medskip

\emph{Note: for each Part (a)-(d), the template provides the corresponding proof or algorithm
  for \prob{Activity Selection} as a starting point.  You can edit that proof or algorithm to create
  your own answer.}

\setcounter{lemma}{0}
\setcounter{theorem}{0}

\continued

\inputAnswer{problem_1a_answer}
\continued

\inputAnswer{problem_1b_answer}
\continued

\inputAnswer{problem_1c_answer}

\vfill

\inputAnswer{problem_1d_answer}

\newpage

%%%%%%%%%%%%%%%%%%%%%%%%%%%%%%%%%%%%%%%%%%%%%%%%%%%%%%%%%%%% 
%%%%%%%%%%%%%%%%%%%%%%%%%%%%%%%%%%%%%%%%%%%%%%%%%%%%%%%%%%%% 
%%%%%%%%%%%%%%%%%%%%%%%%%%%%%%%%%%%%%%%%%%%%%%%%%%%%%%%%%%%% 
%%%%%%%%%%%%%%%%%%%%%%%%%%%%%%%%%%%%%%%%%%%%%%%%%%%%%%%%%%%% 

%%%%% EDIT THE FILES NAMED problem_X_answer.tex.  DON'T EDIT THIS FILE 

\noindent\adjustbox{valign=t}{\parbox{4.6in}{
\problem\label{exercise: BFS}\emph{(Exercise 6.3; BFS)}
Simulate BFS on the graph to the right, starting from the node $a$.
When there are multiple choices for the next node to visit, break ties alphabetically
(choose $b$ before $c$, etc.).
\begin{enumerate}[label=(\alph*)]
% [review 2026-09-27] fix: removed stray "p" at the end of the sentence
\item List the nodes in each level $\ell$ (at distance $\ell$ from $a$).
\item For each node other than $a$, give its parent in the BFS tree.
\item Optionally, draw the BFS tree.
  If you do this, and your answers for Parts (a) and (b)
  are substantially incorrect, you may get partial credit here.
\end{enumerate}
}}\adjustbox{valign=t}{
\begin{tikzpicture}[node distance = 0.6 cm and 1.75cm,on grid,
thick,state/.style ={circle,draw=black,minimum width =.3 cm}]
\node[state] (a) {a};
\node[state] (b) [yshift = 30pt, right =of a] {b};
\node[state] (c) [yshift = -30pt, right =of b] {c};
\node[state] (d) [below =of a, yshift=-14pt] {d};
\node[state] (e) [yshift = 20pt, right =of d] {e};
\node[state] (f) [yshift = -20pt, right =of e] {f};
\node[state] (g) [below =of e, xshift=-27pt,yshift=-25pt] {g};
\node[state] (h) [right =of g] {h};
\path (a) edge [bend left = 0] (b) 
		edge [bend left = 0] (c) 
		edge (d);
\path (b) edge [bend left = 0] (f);
\path (c) edge [bend left = 0] (e);
% [review 2026-09-27] fix: removed the invisible TikZ self-loop "(d) edge (d)" (not an edge of the graph)
\path (d) edge [bend right = 30] (c) 
		edge [bend left = 0] (e);
\path (e) edge [bend left = 0] (h);
\path (g) edge [bend left = 0] (h);
\end{tikzpicture}
}
\refreshHeader 

\setlength{\ITEMHT}{2in}
{\injectEnumerate
  \inputAnswer{problem_2_answer}
}

%%%%%%%%%%%%%%%%%%%%%%%%%%%%%%%%%%%%%%%%%%%%%%%%%%%%%%%%%%%% 
%%%%%%%%%%%%%%%%%%%%%%%%%%%%%%%%%%%%%%%%%%%%%%%%%%%%%%%%%%%% 
%%%%%%%%%%%%%%%%%%%%%%%%%%%%%%%%%%%%%%%%%%%%%%%%%%%%%%%%%%%% 
%%%%%%%%%%%%%%%%%%%%%%%%%%%%%%%%%%%%%%%%%%%%%%%%%%%%%%%%%%%% 

%%%%% EDIT THE FILES NAMED problem_X_answer.tex.  DON'T EDIT THIS FILE 

\newpage 
\noindent\adjustbox{valign=t}{\parbox{4.6in}{\small
    \problem\emph{(Exercise 6.4; DFS)}
    Simulate DFS (as described in the lecture notes) on the graph to the right.
    When there are multiple choices for the next node to visit, break ties alphabetically
    (choose $a$ before $b$, choose $b$ before $c$, etc.).

    \begin{enumerate}[label=(\alph*)]
    \item What is the parent of each node in the resulting DFS forest?
      \\(Write ``none'' for each root.)
      
    \item What's the post-order number (finishing time, Section 6.3.4), of each node?

    \item Which edges are forward edges?  Back edges?  Cross edges? 

    \item Optionally, draw the DFS tree, and label each node with its post-order number.
      If you do this, and your answers for Parts (a)--(c)
      are substantially incorrect, you may get partial credit here.

      % \item Is this graph a directed acyclic graph (DAG)?
    \end{enumerate}
  }}
\adjustbox{valign=t}{
\begin{tikzpicture}[>= latex, node distance = 1 cm and 1.75cm,on grid, thick,state/.style ={circle,draw=black,black,text=black,minimum width =.3 cm}]
\node[state] (a) {a};
\node[state] (b) [right =of a] {b};
\node[state] (c) [right =of b] {c};
\node[state] (d) [below =of a, yshift=-10] {d};
\node[state] (e) [right =of d] {e};
\node[state] (f) [right =of e] {f};
\node[state] (g) [below =of d, yshift=-10] {g};
\node[state] (h) [right =of g] {h};
\node[state] (i) [right =of h] {i};
\tikzset{every node/.style={fill=white}} 
\path[->] 	(a) edge  				(b)
		(a) edge				(e)
		(b) edge				(c)
		(c) edge [bend right=35]	(a)
		(d) edge				(a)
		(d) edge				(e)
		(e) edge				(b)
		(e) edge				(c)
		(e) edge				(f)
		(e) edge				(g)
		(f) edge				(c)
		(g) edge				(d)
		(g) edge				(h)
		(h) edge				(e)
		(h) edge				(i)
		(i) edge				(f);
\end{tikzpicture}
}
% \begin{enumerate}
% \item[(e)]
%   If so, give the topological ordering that the DFS-based algorithm (from class and this note)
%   would output after doing the DFS that you simulated above.

%   If not, how many strongly connected components does the graph have?  List them all.
% \end{enumerate}
\refreshHeader 

\setlength{\ITEMHT}{2.1in}
{\injectEnumerate
\inputAnswer{problem_3_answer}
}
\end{document}

%%% Local Variables:
%%% mode: latex
%%% TeX-master: t
%%% End:
